% ===== PRÉAMBULE CANONIQUE — à coller VERBATIM en tête de CHAQUE .tex =====
\documentclass[11pt,a4paper]{article}
\usepackage[utf8]{inputenc}\usepackage[T1]{fontenc}\usepackage{lmodern}
\usepackage[french]{babel}
\usepackage{amsmath,amssymb,mathtools}\usepackage{icomma}
\usepackage{siunitx}\sisetup{locale=FR}  % \qty{}{}, \si{}, \ang{}
\usepackage{xcolor}\usepackage{colortbl}
\usepackage{tikz}\usepackage{pgfplots}\pgfplotsset{compat=1.18}
\usepackage[siunitx,europeanresistors]{circuitikz} % circuits (élec.)
\usepackage[most]{tcolorbox}\usepackage{enumitem}\usepackage{array,booktabs}
\usepackage[a4paper,margin=2.2cm]{geometry}
\usetikzlibrary{arrows.meta,calc,angles,quotes,positioning,patterns,%
  backgrounds,shapes.geometric,decorations.pathmorphing,decorations.markings,intersections}
\definecolor{primary}{RGB}{222,49,99}\definecolor{primarydark}{RGB}{150,22,55}
\definecolor{defcol}{RGB}{0,114,178}\definecolor{methcol}{RGB}{52,92,148}
\definecolor{excol}{RGB}{0,135,90}\definecolor{attncol}{RGB}{200,40,55}
\tcbset{boxbase/.style={enhanced,breakable,boxrule=0.9pt,arc=2.5pt,
  left=3mm,right=3mm,top=2.3mm,bottom=2.3mm,fonttitle=\bfseries,coltitle=white}}
% --- boîtes TUTORIEL (noms EXACTS, ne pas renommer) ---
\newtcolorbox{cle}[1][]{boxbase,colback=defcol!6,colframe=defcol,
  title={\textsc{À retenir}\ifx\relax#1\relax\relax\else\textmd{ : \itshape #1}\fi}}
\newtcolorbox{methode}[1][]{boxbase,colback=methcol!7,colframe=methcol,sharp corners,
  title={\textsc{Méthode}\ifx\relax#1\relax\relax\else\textmd{ --- \itshape #1}\fi}}
\newtcolorbox{exemplebox}[1][]{boxbase,colback=excol!5,colframe=excol,
  title={\textsc{Exemple}\ifx\relax#1\relax\relax\else\textmd{ : \itshape #1}\fi}}
\newtcolorbox{piege}[1][]{boxbase,colback=attncol!6,colframe=attncol,
  title={\textsc{Piège}\ifx\relax#1\relax\relax\else\textmd{ : \itshape #1}\fi}}
\newtcolorbox{remarque}[1][]{boxbase,colback=black!4,colframe=black!45,
  title={\textsc{Remarque}\ifx\relax#1\relax\relax\else\textmd{ : \itshape #1}\fi}}
% --- boîtes EXERCICE / CORRIGÉ (noms EXACTS, auto-comptées) ---
\newtcolorbox[auto counter]{exo}[1][]{boxbase,colback=primary!4,colframe=primary,
  title={\textsc{Exercice}~\thetcbcounter\ifx\relax#1\relax\relax\else\textmd{ --- \itshape #1}\fi}}
\newtcolorbox[auto counter]{corrige}[1][]{boxbase,colback=excol!4,colframe=excol,
  title={\textsc{Corrigé}~\thetcbcounter\ifx\relax#1\relax\relax\else\textmd{ --- \itshape #1}\fi}}
\newcommand{\vv}[1]{\vec{#1}}\newcommand{\ds}{\displaystyle}
\newcommand{\R}{\mathbb{R}}\newcommand{\N}{\mathbb{N}}\newcommand{\Z}{\mathbb{Z}}
% (un \usepackage{fancyhdr}+en-tête est possible ; il est ignoré côté web)
% =========================================================================
\begin{document}
\begin{corrige}[lire les lignes de champ d'un aimant]
\begin{enumerate}[label=\alph*)]
  \item À l'extérieur, les lignes de champ sont orientées \textbf{du pôle Nord vers le pôle Sud}
        (elles se referment en boucles, repassant du S au N \emph{à l'intérieur} de l'aimant).
  \item Au point $A$ : c'est là que le champ est le \textbf{plus intense}. Près des pôles, les lignes
        de champ sont \textbf{plus resserrées} ; or plus les lignes sont serrées, plus $\vv{B}$ est
        grand. Au point $C$, éloigné, les lignes s'écartent : le champ y est faible.
\end{enumerate}
\end{corrige}

\begin{corrige}[sens de $\vv{B}$ autour d'un fil : main droite]
On applique la règle de la main droite : \textbf{pouce vers le haut} (sens de $I$), les doigts
s'enroulent autour du fil.
\begin{enumerate}[label=\alph*)]
  \item Du côté droit, les doigts « plongent » derrière le fil : au point $M$, le champ
        \textbf{entre} dans la feuille, $\otimes$.
  \item Du côté gauche, les doigts « remontent » vers l'avant : au point $N$, le champ \textbf{sort}
        de la feuille, $\odot$. (Les deux sont bien de sens opposés, comme attendu de part et
        d'autre du fil.)
\end{enumerate}
\end{corrige}

\begin{corrige}[superposition de deux champs alignés]
La superposition est une somme \textbf{vectorielle}. Les deux champs étant colinéaires :
\begin{enumerate}[label=\alph*)]
  \item Même sens : les normes s'ajoutent,
        \[ B_{\text{tot}}=B_1+B_2=3\cdot10^{-5}+2\cdot10^{-5}=\qty{5e-5}{\tesla}. \]
  \item Sens opposés : les normes se \textbf{soustraient},
        \[ B_{\text{tot}}=|B_1-B_2|=|3\cdot10^{-5}-2\cdot10^{-5}|=\qty{1e-5}{\tesla}, \]
        dirigé dans le sens du \textbf{plus grand}, c'est-à-dire celui de $\vv{B}_1$.
\end{enumerate}
\end{corrige}

\begin{corrige}[sens de $\vv{B}$ dans une spire]
Règle de la main droite « version spire » : on courbe les \textbf{doigts dans le sens de $I$}
(ici trigonométrique) ; le \textbf{pouce} donne le sens de $\vv{B}$ à l'intérieur.
\begin{enumerate}[label=\alph*)]
  \item Les doigts tournant dans le sens trigonométrique, le pouce pointe \textbf{vers toi} :
        $\vv{B}$ \textbf{sort} de la feuille, $\odot$.
  \item Le champ sort toujours par la face \textbf{Nord} : c'est donc la face \textbf{tournée vers toi}
        qui est la face Nord (la face de derrière est Sud).
\end{enumerate}
\end{corrige}

\begin{corrige}[fil vu en bout : convention $\odot/\otimes$]
Le courant sort de la feuille ($\odot$) : on place le \textbf{pouce vers soi}, les doigts s'enroulent.
\begin{enumerate}[label=\alph*)]
  \item Vus de face avec le pouce pointé vers soi, les doigts tournent dans le sens
        \textbf{trigonométrique} (anti-horaire).
  \item Sur le cercle parcouru dans le sens anti-horaire, la tangente au point le plus haut $P$
        pointe vers la \textbf{gauche} : $\vv{B}$ est dirigé \textbf{vers la gauche} en $P$.
\end{enumerate}
\end{corrige}

\end{document}
